\subsection{AI Asisstant Concept}

    Although blockchain technology provides a secure and immutable mechanism for certificate verification, the underlying technical information may be difficult for non-technical users to understand. Verification results often contain transaction hashes, wallet addresses, token identifiers, and IPFS references that are unfamiliar to recruiters, employers, or university administrators.\\

    To improve usability, this system introduces an AI-Assisted Certificate Verification Module. The purpose of this module is not to perform blockchain verification independently, but rather to translate verified blockchain data into clear and understandable natural language explanations.\\

    The AI assistant acts as an interpretation layer between the blockchain verification system and end users. After certificate information has been validated by the application, the assistant analyzes the verified data and provides human-readable explanations regarding:

    \begin{itemize}
        \item Certificate authenticity status.
        \item NFT ownership information.
        \item Blockchain transaction records.
        \item IPFS metadata storage references.
        \item Academic certificate details.
        \item Certificate lifecycle status.
    \end{itemize}

    This approach allows users to interact with blockchain verification results through conversational questions rather than manually interpreting technical records.
